% =============================================================
% ESTILO VISUAL GENERAL
% INFORMES DE LABORATORIO ANSYS / TELECOMUNICACIONES
% =============================================================

% -------------------------------------------------------------
% PALETA DE COLORES
% -------------------------------------------------------------

\definecolor{Primary}{HTML}{0E7490}
\definecolor{PrimaryDark}{HTML}{134E4A}

\definecolor{Accent}{HTML}{ECFEFF}
\definecolor{SoftGray}{HTML}{F8FAFC}

\definecolor{TextGray}{HTML}{334155}
\definecolor{RuleGray}{HTML}{CBD5E1}
\definecolor{MutedGray}{HTML}{64748B}

% Colores semánticos para tarjetas de diseño moderno
\definecolor{WarnBorder}{HTML}{D97706} % Amber/Oro
\definecolor{WarnBG}{HTML}{FFFBEB}     % Fondo Oro ultra-suave
\definecolor{SuccessBorder}{HTML}{059669} % Esmeralda/Verde
\definecolor{SuccessBG}{HTML}{F0FDF4}     % Fondo Verde ultra-suave
\definecolor{InfoBorder}{HTML}{4F46E5}    % Indigo/Azul
\definecolor{InfoBG}{HTML}{F5F3FF}        % Fondo Indigo ultra-suave

% -------------------------------------------------------------
% PDF
% -------------------------------------------------------------

\hypersetup{
    colorlinks=true,
    linkcolor=PrimaryDark,
    citecolor=PrimaryDark,
    urlcolor=PrimaryDark,
    pdftitle={\TituloInforme},
    hypertexnames=false
}
% =============================================================
% TEXTO GENERAL
% =============================================================

\onehalfspacing

\setlength{\parindent}{0pt}
\setlength{\parskip}{0.40em}

% =============================================================
% LISTAS
% =============================================================

\setlist[itemize]{
    topsep=3pt,
    itemsep=2pt,
    parsep=0pt,
    partopsep=0pt,
    leftmargin=1.15cm
}

\setlist[enumerate]{
    topsep=3pt,
    itemsep=2pt,
    parsep=0pt,
    partopsep=0pt,
    leftmargin=1.15cm
}

% =============================================================
% TÍTULOS
% =============================================================

\titleformat{\section}
{\Large\bfseries\color{Primary}}
{\thesection.}
{0.75em}
{}
[
\vspace{0.16em}
{\color{Primary}\titlerule[0.6pt]}
]

\titleformat{\subsection}
{\large\bfseries\color{PrimaryDark}}
{\thesubsection.}
{0.60em}
{}

\titleformat{\subsubsection}
{\normalsize\bfseries\color{TextGray}}
{\thesubsubsection.}
{0.50em}
{}

\titlespacing*{\section}
{0pt}
{1.70em}
{0.75em}

\titlespacing*{\subsection}
{0pt}
{1.00em}
{0.40em}

\titlespacing*{\subsubsection}
{0pt}
{0.70em}
{0.25em}

% =============================================================
% ENCABEZADO Y PIE
% =============================================================

\fancypagestyle{informe}{
    \fancyhf{}
    \fancyfoot[C]{\scriptsize\color{MutedGray}\thepage\ de\ \pageref*{LastPage}}
    \renewcommand{\headrulewidth}{0pt}
    \renewcommand{\footrulewidth}{0pt}
}

\pagestyle{informe}

\renewcommand{\sectionmark}[1]{%
    \markboth{\thesection.\ #1}{\thesection.\ #1}
}
\renewcommand{\subsectionmark}[1]{}

% Para índice y páginas sin encabezado (TOC, Referencias, etc.)
\fancypagestyle{plain}{
    \fancyhf{}
    \fancyfoot[C]{\scriptsize\color{MutedGray}\thepage}
    \renewcommand{\headrulewidth}{0pt}
    \renewcommand{\footrulewidth}{0pt}
}

% =============================================================
% CONFIGURACIÓN DE CAPTIONS (TÍTULOS DE FIGURAS Y TABLAS)
% =============================================================

\captionsetup{
    font=small,
    labelfont={bf,color=PrimaryDark},
    textfont=it,
    justification=centering,
    singlelinecheck=false
}
\captionsetup[figure]{aboveskip=8pt, belowskip=4pt}
\captionsetup[table]{aboveskip=6pt, belowskip=8pt} % Espacio ordenado debajo del título de la tabla

% =============================================================
% COMPORTAMIENTO Y ESPACIADO DE ELEMENTOS FLOTANTES (FIGURAS/TABLAS)
% =============================================================

% Evita que las figuras se acumulen al final o dejen páginas vacías
\renewcommand{\topfraction}{0.85}
\renewcommand{\bottomfraction}{0.85}
\renewcommand{\textfraction}{0.15}
\renewcommand{\floatpagefraction}{0.75}

% Distancia ordenada entre flotantes y el texto
\setlength{\floatsep}{12pt plus 2pt minus 2pt}
\setlength{\textfloatsep}{14pt plus 2pt minus 4pt}
\setlength{\intextsep}{12pt plus 2pt minus 2pt}

% Espaciado compacto para ecuaciones (evita grandes huecos blancos en interlineado 1.5)
\AtBeginDocument{
    \setlength{\abovedisplayskip}{6pt plus 2pt minus 2pt}
    \setlength{\belowdisplayskip}{6pt plus 2pt minus 2pt}
    \setlength{\abovedisplayshortskip}{0pt plus 2pt}
    \setlength{\belowdisplayshortskip}{4pt plus 2pt minus 2pt}
}

% =============================================================
% TABLAS
% =============================================================

\renewcommand{\arraystretch}{1.15}
\setlength{\tabcolsep}{7pt}

% =============================================================
% CAJAS
% =============================================================

\tcbset{
    colback=Accent,
    colframe=Primary,
    boxrule=0.4pt,
    arc=1.5mm,
    left=6pt,
    right=6pt,
    top=5pt,
    bottom=5pt
}

% =============================================================
% CÓDIGO
% =============================================================

\newcommand{\noncopynumber}[1]{%
    \BeginAccSupp{ActualText={}}%
    \tiny\color{MutedGray}#1%
    \EndAccSupp{}%
}

\lstdefinestyle{labcode}{
    backgroundcolor=\color{SoftGray},
    basicstyle=\ttfamily\small,
    breaklines=true,
    columns=flexible,
    keepspaces=true,
    frame=none,
    rulecolor=\color{RuleGray},
    keywordstyle=\bfseries\color{PrimaryDark},
    commentstyle=\itshape\color{MutedGray},
    stringstyle=\color{Primary},
    numbers=left,
    numberstyle=\noncopynumber,
    xleftmargin=1.3em,
    framexleftmargin=1.0em,
    aboveskip=0.6em,
    belowskip=0.6em
}

\lstset{style=labcode}

% =============================================================
% CONFIGURACIÓN DE ÍNDICES (TOC, LOF, LOT)
% =============================================================

% Puntos de guía (dot leaders) para secciones, subsecciones y subsubsecciones
\renewcommand{\cftsecleader}{\cftdotfill{\cftdotsep}}
\renewcommand{\cftsubsecleader}{\normalfont\cftdotfill{\cftdotsep}}
\renewcommand{\cftsubsubsecleader}{\normalfont\cftdotfill{\cftdotsep}}

% Estilo de fuentes para Sección (negrita y color principal)
\renewcommand{\cftsecfont}{\bfseries\color{PrimaryDark}}
\renewcommand{\cftsecpagefont}{\bfseries\color{PrimaryDark}}

% Estilo de fuentes para Subsección y Subsubsección
\renewcommand{\cftsubsecfont}{\normalfont\color{TextGray}}
\renewcommand{\cftsubsecpagefont}{\normalfont\color{TextGray}}
\renewcommand{\cftsubsubsecfont}{\normalfont\color{MutedGray}}
\renewcommand{\cftsubsubsecpagefont}{\normalfont\color{MutedGray}}

% Ajuste de espaciado vertical
\setlength{\cftbeforesecskip}{0.6em}
\setlength{\cftbeforesubsecskip}{0.25em}
\setlength{\cftbeforesubsubsecskip}{0.15em}

% Configuración para Índice de Figuras y Tablas
\renewcommand{\cftfigfont}{\normalfont\color{TextGray}}
\renewcommand{\cftfigpagefont}{\normalfont\color{TextGray}}
\renewcommand{\cftfigleader}{\normalfont\cftdotfill{\cftdotsep}}

\renewcommand{\cfttabfont}{\normalfont\color{TextGray}}
\renewcommand{\cfttabpagefont}{\normalfont\color{TextGray}}
\renewcommand{\cfttableader}{\normalfont\cftdotfill{\cftdotsep}}

% Convertir secciones a mayúsculas automáticamente en el índice (TOC)
\let\oldaddcontentsline\addcontentsline
\renewcommand{\addcontentsline}[3]{%
    \ifstrequal{#2}{section}%
        {\oldaddcontentsline{#1}{#2}{\MakeUppercase{#3}}}%
        {\oldaddcontentsline{#1}{#2}{#3}}%
}



